\documentclass[10pt,a4paper]{article}

\usepackage[T1]{fontenc}
\usepackage[utf8]{inputenc}
\usepackage[ngerman]{babel}
\usepackage[a4paper,margin=12mm]{geometry}
\usepackage{lmodern}
\usepackage[table]{xcolor}
\usepackage{tabularx}
\usepackage{array}
\usepackage{booktabs}
\usepackage{amsmath}
\usepackage{microtype}
\usepackage[most]{tcolorbox}
\usepackage{tikz}
\usetikzlibrary{arrows.meta,calc,positioning}

\pagestyle{empty}
\setlength{\parindent}{0pt}
\setlength{\parskip}{0pt}
\renewcommand{\familydefault}{\sfdefault}
\renewcommand{\arraystretch}{1.05}

\definecolor{Navy}{HTML}{17324D}
\definecolor{Blue}{HTML}{3B6EA5}
\definecolor{Sky}{HTML}{EAF3FA}
\definecolor{Mint}{HTML}{E8F5EF}
\definecolor{Gold}{HTML}{F3B33D}
\definecolor{Ink}{HTML}{1E2935}
\definecolor{Muted}{HTML}{617080}
\definecolor{Line}{HTML}{D8E1E8}

\color{Ink}
\newcommand{\sectiontitle}[1]{%
  \vspace{0.4mm}{\color{Blue}\bfseries\large #1}\par
  \vspace{0.15mm}{\color{Blue}\rule{\linewidth}{0.7pt}}\par
  \vspace{0.2mm}}
\newcommand{\unit}[1]{\,\mathrm{#1}}
\newcommand{\unitfrac}[2]{\,\frac{\mathrm{#1}}{\mathrm{#2}}}
\newcommand{\pagefooter}[1]{%
  \vfill{\color{Line}\rule{\linewidth}{0.5pt}}\par\vspace{1mm}
  \begin{tabularx}{\linewidth}{@{}Xc r@{}}
    {\leavevmode\small\color{Muted}Klasse 11 · Physik · Kreisbewegungen} &
    {\leavevmode\small\color{Muted}#1} &
    {\leavevmode\small\color{Muted}Sicherungsblatt · Aufgabe 4}
  \end{tabularx}}

\begin{document}

\colorbox{Navy}{%
  \parbox{\dimexpr\linewidth-2\fboxsep\relax}{%
    \vspace{1.3mm}\color{white}
    \begin{tabularx}{\linewidth}{@{}Xr@{}}
      {\small\bfseries PHYSIK · KLASSE 11} &
      \colorbox{Gold}{\color{Navy}\small\bfseries\strut LÖSUNGEN}
    \end{tabularx}
    \vspace{0.9mm}
    {\fontsize{18}{21}\selectfont\bfseries Aufgabe 4 -- Zentripetalkraft}\par
    \vspace{0.3mm}{\large Sicherungsblatt}\vspace{1.1mm}
  }%
}

\vspace{0.9mm}
\colorbox{Sky}{\parbox{\dimexpr\linewidth-2\fboxsep\relax}{%
  \vspace{0.55mm}\textcolor{Blue}{\bfseries Ziel:} Die Richtung der Zentripetalkraft kennen und ihren Betrag bestimmen.\vspace{0.55mm}}}

\sectiontitle{1 · Einflussgrößen}
\small
\rowcolors{2}{Sky}{white}
\begin{tabularx}{\linewidth}{@{}>{\bfseries}p{27mm} X >{\centering\arraybackslash}p{25mm}@{}}
\toprule[0.4pt]
Größe & Je-desto-Aussage & Proportionalität \\
\midrule[0.3pt]
$m$ (r, $\omega$ gleich) & Je größer $m$, desto größer $F_{\mathrm Z}$. & $F_{\mathrm Z}\sim m$ \\
$\omega$ (m, r gleich) & Je größer $\omega$, desto größer $F_{\mathrm Z}$. Doppeltes $\omega$ $\to$ vierfache Kraft. & $F_{\mathrm Z}\sim\omega^{2}$ \\
$r$ bei gleichem $\omega$ & Je größer $r$, desto größer $F_{\mathrm Z}$. & $F_{\mathrm Z}\sim r$ \\
$r$ bei gleichem $v_{\mathrm B}$ & Je größer $r$, desto kleiner $F_{\mathrm Z}$. & $F_{\mathrm Z}\sim\dfrac{1}{r}$ \\
\bottomrule[0.4pt]
\end{tabularx}
\vspace{1.5mm}

\sectiontitle{2 · Merksatz}
\begin{tcolorbox}[colback=Sky,colframe=Line,arc=1.2mm,boxrule=.6pt,left=2mm,right=2mm,top=.7mm,bottom=.7mm]
{\centering\fontsize{14}{17}\selectfont $F_{\mathrm Z}=\dfrac{m\cdot v_{\mathrm B}^{2}}{r}=m\cdot\omega^{2}\cdot r$\par}
\vspace{0.3mm}
{\centering\small $m$: Masse (kg) \quad $v_{\mathrm B}$: Bahngeschwindigkeit $\left(\unitfrac{m}{s}\right)$ \quad $\omega$: Winkelgeschwindigkeit $\left(\unitfrac{rad}{s}\right)$ \quad $r$: Radius (m)\par}
\vspace{0.2mm}
\small Bewegt sich ein Körper der Masse $m$ mit der Bahngeschwindigkeit $v_{\mathrm B}$ auf einer Kreisbahn mit dem Radius $r$, hat die nötige Zentripetalkraft den Betrag $F_{\mathrm Z}=\dfrac{m\cdot v_{\mathrm B}^{2}}{r}=m\cdot\omega^{2}\cdot r$.
\end{tcolorbox}

\sectiontitle{3 · Richtung der Zentripetalkraft}
\begin{tcolorbox}[colback=Sky,colframe=Line,arc=1.2mm,boxrule=.6pt,left=2mm,right=2mm,top=.7mm,bottom=.7mm]
\small Damit sich ein Körper mit gleichbleibender Bahngeschwindigkeit auf einer Kreisbahn bewegt, muss auf ihn eine Kraft wirken, die zum \textbf{Mittelpunkt} hin gerichtet ist. Diese Kraft nennt man Zentripetalkraft $F_{\mathrm Z}$.
\end{tcolorbox}
\vspace{1mm}
\colorbox{Gold}{\parbox{\dimexpr\linewidth-2\fboxsep\relax}{\small\centering\bfseries Fehlt die Kraft zum Mittelpunkt $\to$ tangential geradlinig weiter · Kraft zu klein $\to$ Bahnradius wird größer}}

\vspace{1.5mm}
\sectiontitle{4 · Rechenbeispiele}
{\renewcommand{\arraystretch}{1.42}%
\begin{tcbraster}[raster columns=2,raster equal height,raster column skip=3mm]
\begin{tcolorbox}[colback=white,colframe=Line,arc=1.2mm,boxrule=.65pt,left=2.2mm,right=2.2mm,top=1.1mm,bottom=1.1mm]
\textcolor{Blue}{\bfseries Hammerwurf}\par\vspace{0.35mm}
{\small
\begin{tabularx}{\linewidth}{@{}>{\bfseries}p{19mm}X@{}}
Gegeben: & $m=7{,}26\unit{kg}$, $v_{\mathrm B}=25\unitfrac{m}{s}$, $r=1{,}8\unit{m}$ \\
Gesucht: & $F_{\mathrm Z}$ \\
Ansatz: & $F_{\mathrm Z}=\dfrac{m\cdot v_{\mathrm B}^{2}}{r}$ \\[2mm]
Rechnung: & $F_{\mathrm Z}=\dfrac{7{,}26\unit{kg}\cdot(25\unitfrac{m}{s})^{2}}{1{,}8\unit{m}}$
\end{tabularx}}
\par\vspace{1.2mm}\hfill\textcolor{Blue}{\bfseries Ergebnis: $2\,521\unit{N}\approx 2{,}5\unit{kN}$ \quad 2 gültige Ziffern}
\end{tcolorbox}
\begin{tcolorbox}[colback=white,colframe=Line,arc=1.2mm,boxrule=.65pt,left=2.2mm,right=2.2mm,top=1.1mm,bottom=1.1mm]
\textcolor{Blue}{\bfseries Auto in der Kurve}\par\vspace{0.35mm}
{\small
\begin{tabularx}{\linewidth}{@{}>{\bfseries}p{19mm}X@{}}
Gegeben: & $m=1200\unit{kg}$, $54\unitfrac{km}{h}=15\unitfrac{m}{s}$, $r=40\unit{m}$ \\
Gesucht: & $F_{\mathrm Z}$ \\
Ansatz: & $F_{\mathrm Z}=\dfrac{m\cdot v_{\mathrm B}^{2}}{r}$ \\[2mm]
Rechnung: & $F_{\mathrm Z}=\dfrac{1200\unit{kg}\cdot(15\unitfrac{m}{s})^{2}}{40\unit{m}}$
\end{tabularx}}
\par\vspace{1.2mm}\hfill\textcolor{Blue}{\bfseries Ergebnis: $6\,750\unit{N}\approx 6{,}8\unit{kN}$ \quad 2 gültige Ziffern}
\end{tcolorbox}
\end{tcbraster}

\vspace{3mm}
\begin{tcolorbox}[colback=white,colframe=Line,arc=1.2mm,boxrule=.65pt,left=2.2mm,right=2.2mm,top=1.1mm,bottom=1.1mm]
\textcolor{Blue}{\bfseries Höchstgeschwindigkeit}\par\vspace{0.35mm}
{\small
\begin{tabularx}{\linewidth}{@{}>{\bfseries}p{19mm}X@{}}
Gegeben: & $m=1200\unit{kg}$, $r=40\unit{m}$, $F_{\mathrm Z}=9{,}6\unit{kN}=9\,600\unit{N}$ \\
Gesucht: & $v_{\mathrm B}$ \\
Ansatz: & $v_{\mathrm B}=\sqrt{\dfrac{F_{\mathrm Z}\cdot r}{m}}$ \\[2mm]
Rechnung: & $v_{\mathrm B}=\sqrt{\dfrac{9\,600\unit{N}\cdot 40\unit{m}}{1200\unit{kg}}}=\sqrt{320}\unitfrac{m}{s}$
\end{tabularx}}
\par\vspace{1.2mm}\hfill\textcolor{Blue}{\bfseries Ergebnis: $\approx 18\unitfrac{m}{s}$ \quad ($17{,}89\unitfrac{m}{s}\cdot 3{,}6\approx 64\unitfrac{km}{h}$) \quad 2 gültige Ziffern}
\end{tcolorbox}}

\pagefooter{1 / 1}

\end{document}
